\Askhsh{20678}{4}
\begin{ylh}{1}
    \item 8.2. Κριτήρια ομοιότητας
    \item 10.4. Άλλοι τύποι για το εμβαδόν τριγώνου
    \item 10.5. Λόγος εμβαδών όμοιων τριγώνων - πολυγώνων.
\end{ylh}
\begin{wrapfigure}[15]{r}{5.5cm}
    \centering
    \begin{tikzpicture}[scale=1, every node/.style={font=\small}]
        % Define center
        \tkzDefPoint(0,0){O}
        % Define outer rectangle vertices
        \tkzDefPoint(-2.5,1.8){A}
        \tkzDefPoint(2.5,1.8){B}
        \tkzDefPoint(2.5,-1.8){C} % \Gamma
        \tkzDefPoint(-2.5,-1.8){D} % \Delta
        
        % Define inner rectangle vertices via homothety (ratio approx 1/sqrt(2) ~= 0.707)
        \tkzDefPointBy[homothety=center O ratio 0.7](A){A'}
        \tkzDefPointBy[homothety=center O ratio 0.7](B){B'}
        \tkzDefPointBy[homothety=center O ratio 0.7](C){C'}
        \tkzDefPointBy[homothety=center O ratio 0.7](D){D'}
        
        % Draw elements
        \tkzDrawPolygon(A,B,C,D)
        \tkzDrawPolygon(A',B',C',D')
        \tkzDrawSegments[dashed](A,C B,D)
        
        % Draw points
        \tkzDrawPoints(O,A,B,C,D,A',B',C',D')

        % Label points
        \tkzLabelPoint[above left](A){A}
        \tkzLabelPoint[above right](B){B}
        \tkzLabelPoint[below right](C){$\Gamma$}
        \tkzLabelPoint[below left](D){$\Delta$}
        \tkzLabelPoint[above left, yshift=-1mm, xshift=2mm](A'){A'}
        \tkzLabelPoint[above right, yshift=-1mm, xshift=-2mm](B'){B'}
        \tkzLabelPoint[below right, yshift=1.5mm, xshift=-2mm](C'){$\Gamma'$}
        \tkzLabelPoint[below left, yshift=1.5mm, xshift=2mm](D'){$\Delta'$}
        \tkzLabelPoint[above right, xshift=-2mm](O){O}
    \end{tikzpicture}
\end{wrapfigure}
Η κορνίζα του παρακάτω σχήματος αποτελείται από δύο όμοια ορθογώνια με παράλληλες πλευρές και κοινό κέντρο O. Το ορθογώνιο Α'Β'Γ'Δ' έχει το μισό εμβαδόν από το ορθογώνιο ΑΒΓΔ.
\begin{alist}
    \item Να βρείτε τον λόγο ομοιότητας του ορθογωνίου ΑΒΓΔ προς το ορθογώνιο Α'Β'Γ'Δ'. \monades{5}
    \item Να αποδείξετε ότι τα τρίγωνα ΑΒΓ και Α'Β'Γ' είναι όμοια. \monades{6}
    \item Στην κορνίζα τοποθετούμε μια φωτογραφία που χωράει ακριβώς στο κάδρο, χωρίς να χάνεται κανένα μέρος της. Η διαγώνιος ΑΓ της κορνίζας έχει μήκος 40 cm και $\widehat{ΑΟΒ} = 120^\circ$.
    \begin{rlist}
        \item Πόσο μήκος έχει η διαγώνιος της φωτογραφίας; \monades{6}
        \item Πόσο είναι το εμβαδόν της φωτογραφίας; \monades{8}
    \end{rlist}
\end{alist}